
lemme facile suivant, donc pour @@M0@@ en vertu de 3.5 et 3.4 (ii), enfin pour @@M1@@ gr\^ace \`a 3.4 (i).

\textbf{Lemme 3.7.} \emph{Soient @@M2@@ de type fini sur @@M3@@ noeth\'erien et @@M4@@ une famille de classes de faisceaux coh\'erents sur les fibres de @@M5@@. Pour que @@M6@@ soit limit\'ee, il faut et il suffit que pour tout @@M7@@ la famille @@M8@@ induite sur les fibres de @@M9@@ soit limit\'ee.}

Dans la d\'emonstration de 3.6, on s'est ramen\'e au cas o\`u @@M10@@ est projective et int\`egre sur @@M11@@\,; munissons @@M12@@ d'un faisceau tr\`es ample @@M13@@. Soient @@M14@@ une famille limit\'ee de classes de faisceaux inversibles sur les fibres de @@M15@@ et @@M16@@ un faisceau inversible sur un @@M17@@ avec @@M18@@ de type fini sur @@M19@@, qui limite @@M20@@. Quitte \`a couper @@M21@@ en morceaux, on peut supposer @@M22@@ plat sur @@M23@@, et quitte \`a remplacer @@M24@@ par chacune de ses composantes connexes, on peut supposer @@M25@@ connexe. Alors le polyn\^ome en @@M26@@ et @@M27@@, @@M28@@ ne d\'epend pas de @@M29@@ (EGA III 7.9.4).

Soit @@M30@@~: @@M31@@ est la famille des classes des faisceaux inversibles @@M32@@ tels que la classe de @@M33@@ soit dans @@M34@@. Soit @@M35@@, o\`u @@M36@@ est un polyn\^ome en @@M37@@ de degr\'e @@M38@@. Pour @@M39@@ variables, on a @@M40@@ et chaque @@M41@@ ne d\'epend donc pas du @@M42@@ choisi\,; par suite chaque polyn\^ome @@M43@@ ne d\'epend pas non plus de @@M44@@. Par cons\'equent la famille @@M45@@ a un seul polyn\^ome de Hilbert @@M46@@, donc est limit\'ee en vertu de 2.11.

\textbf{Th\'eor\`eme 3.8.} \emph{Soient @@M47@@ un sch\'ema noeth\'erien, @@M48@@ un @@M49@@-sch\'ema projectif muni d'un faisceau inversible relativement ample @@M50@@, @@M51@@ le sch\'ema des z\'eros d'une section @@M52@@ de @@M53@@, et @@M54@@ le morphisme d'inclusion. On}
